\documentclass[11pt]{article}
\usepackage[margin=1in]{geometry}
\usepackage{hyperref}
\usepackage{enumitem}
\title{Temporal Connectivity in Landscape Genetics: From Static Graphs to Time-Series Genomics\\
\large A Literature Review Generated with AI Assistance}
\author{LitReview AI \and A. L. Robles Fernández\\[4pt]
\small LitReview \textbar{} ai.litreview.org\\
\small Identifier: 2608.00003}
\date{2026-08-25}

\begin{document}
\maketitle

\begin{abstract}
Landscape genetics has traditionally treated connectivity as a static property inferred from a single snapshot of genetic variation. Temporal genomic datasets — historical specimens, resampled populations, and time-series of allele frequencies — now allow connectivity itself to be estimated as a dynamic process. This review synthesizes methods linking landscape resistance surfaces to temporal genetic differentiation, including time-series F_ST approaches, effective migration surface models fit across decades, and simulation-based inference for resistance parameterization. We discuss how temporal designs disentangle demographic change from landscape change, quantify lag effects in habitat-loss responses, and forecast gene flow under future land-use scenarios.

\noindent\textbf{Keywords:} landscape genetics, gene flow, temporal genomics, connectivity.
\end{abstract}

\section*{Placeholder Notice}
This document is a placeholder entry in the LitReview repository. The full
review text will be generated with AI assistance from primary literature and
released after curation. This page establishes the identifier, metadata, and
distribution formats (LaTeX, PDF, HTML) for the entry.

\section*{Scope}
The forthcoming review will synthesize recent literature on the topic above,
covering its principal methods, findings, open challenges, and directions for
future work within the field of Ecology \& Evolution.

\section*{Data and Code Availability}
No new data or code are associated with this placeholder.

\bibliographystyle{plain}
\begin{thebibliography}{9}
\bibitem{litreview2026} LitReview Consortium (2026). \emph{LitReview: an arXiv-style repository of AI-generated literature reviews for the basic sciences.} ai.litreview.org.
\end{thebibliography}

\end{document}
